% Source: https://artofproblemsolving.com/community/c3h1567535_regular_pentagon_question
% Stable identity: aops-topic-1567535-q01
In triangle $VWX$, the points $A$ and $B$ are the midpoints of sides $VW$ and $WX$.

\begin{asy}
usepackage("amsmath");
size(5cm);
defaultpen(fontsize(10pt));
pair[] Pp;string[] nm={"V","W","X","Y","Z"};string[] mn={"A","B","C","D","E"};
for(int i=0;i<5;++i) Pp.push(dir(90-72*i));
pair[] M;for(int i=0;i<5;++i) M.push((Pp[i]+Pp[(i+1)%5])/2);
draw(Pp[0]--Pp[1]--Pp[2]--Pp[3]--Pp[4]--cycle,linewidth(1));
draw(M[0]--M[1]--M[2]--M[3]--M[4]--cycle,blue+linewidth(1));
draw(Pp[0]--Pp[2],red+dashed);
for(int i=0;i<5;++i){label("$"+nm[i]+"$",Pp[i],unit(Pp[i]));dot(M[i],blue);label("$"+mn[i]+"$",M[i],unit(M[i]),blue);}
label("$4$",(M[0]+M[1])/2,W,blue);
label("\mbox{}",(min(currentpicture,user=true)+max(currentpicture,user=true))/2,black); // keeps the text after the figure black when LaTeX typesets the labels inline
\end{asy}

By the midline theorem, the segment $AB$ joining them is parallel to the third side $VX$ and half as long (triangle $WAB$ is the image of triangle $WVX$ under the dilation centred at $W$ with ratio $\dfrac{1}{2}$). Therefore
\[
VX=2\cdot AB=2\cdot4=\boxed{8}.
\]
No golden-ratio computations are needed. (If instead the outer pentagon $VWXYZ$ had side $4$, its diagonal would be $4\cdot\dfrac{1+\sqrt5}{2}=2+2\sqrt5$, and the inner pentagon would be irrelevant.)
